
Lorsque @@M0@@ est de caractéristique nulle, i.e. son faisceau structural est un faisceau de @@M1@@-algèbres, le problème se résout simplement en prenant
@@M2@@
où l'on considère la puissance tensorielle @@M3@@.ème (au sens des catégories dérivées) de @@M4@@ comme étant un objet de @@M5@@, @@M6@@ étant l'algèbre du groupe symétrique @@M7@@ à coefficients dans @@M8@@, et où l'exposant alt désigne le passage à la ``partie alternée sous l'action de @@M9@@'' (opération qui a un sens grâce à l'hypothèse de caractéristique nulle). Une formule de ce type était utilisée déjà dans [8] (cité par la suite [RRR]), pour la première démonstration du théorème de Riemann--Roch.

Sans hypothèse de caractéristique, lorsque @@M10@@ est nul en degrés @@M11@@, alors par Dold--Puppe [6] @@M12@@ est canoniquement homotope au complexe associé à un Module semi-simplicial @@M13@@ canoniquement associé à @@M14@@, et un candidat assez naturel pour @@M15@@ (pour @@M16@@ plat, disons) serait le complexe associé au Module semi-simplicial @@M17@@ (dont la composante de degré @@M18@@ est @@M19@@). Mais pour étendre la définition à des complexes de @@M20@@ qui ne sont plus supposés nuls en degrés @@M21@@, il faudrait étudier le comportement de l'opération précédente par rapport à l'opération de translation @@M22@@ (ou \uline{suspension}, dans la terminologie de [6]); malheureusement, il ne semble pas qu'on ait un isomorphisme @@M23@@ qui permettrait de faire de façon triviale l'extension désirée.

\section*{2. La formule de Riemann--Roch sans hypothèses projectives}

\textbf{2.1.} Une fois en possession d'une @@M24@@-structure naturelle sur les @@M25@@, il est possible de formuler le théorème de Riemann--Roch, sous la forme envisagée dans Exp. VIII, pour n'importe quel morphisme @@M26@@ propre et d'intersection complète, (et même, plus généralement, dans le cas d'un schéma relatif @@M27@@ sur un topos localement annelé @@M28@@, @@M29@@ étant encore soumis aux conditions relatives correspondantes). La classe de Todd de @@M30@@ se définit grâce à la classe dans @@M31@@ du complexe cotangent relatif @@M32@@ pour @@M33@@, construit dans Exp. 0 4.4, lequel complexe est ici parfait grâce à l'hypothèse locale faite sur @@M34@@. Il semble clair que la démonstration de la formule de Riemann--Roch dans ce cas général demandera l'introduction d'idées essentiellement nouvelles, et il est possible que ces mêmes idées pourraient donner la clef 
